\section{Limitations}
\label{sec:limitations}

Seventeen limitations, in the order a sceptical reader should apply them, are stated in full in
Appendix~\ref{app:limitations}; the six that most constrain what may be concluded are here.

\textbf{The scale result is within one family.} The $0\to4\to21\to36\to80$ ladder is
Qwen2.5-Coder alone, against one mismatched interface, on one task (item~3). We claim a
within-family monotone relation between checkpoint size and absolute undercount; we do not claim
censoring is generally scale-increasing, and we do not equate parameter count with capability. A
second \emph{family}'s ladder is the test that would promote or refute it.

\textbf{The repaired-FC control is one seed at one budget} (item~8). It is a clean parser-registry
intervention at 7B, but both arms use LoRA and 150 steps on one code task. It establishes the
experience-distribution effect in this setting; its null learning result is not a theorem that a
working channel can never help, nor is it strictly comparable to the historical 1.5B
full-parameter runs.

\textbf{Near-zero samples are not zero policy probability} (items~10--11). The broken 7B arm
retains full rollout text and distinguishes 8{,}689 tight-classified call emissions from 10 accepted and
executed calls, so attempts and discards are directly separated there. The older 1.5B run did not
retain full text and remains over-determined; neither run identifies the support of the policy
beyond its sampled trajectories.

\textbf{Conditioning on a successful parse is post-treatment conditioning} (item~9). Whether an
item parses is determined by the interface under test, so the both-parsed subset is not a random
subsample. On Qwen the residual is $+8.4$\,pp, 95\% CI $[-0.5,+17.4]$, $p$=0.118: \textbf{we do
not detect a residual difference, which is not the same as establishing there is none.}

\textbf{Multiple comparisons} (item~6). Roughly a dozen McNemar tests; the Bonferroni threshold is
$\approx$0.0042. Surviving it: the official-template \texttt{strict: true} intervention on
wrong-tool labels ($p=4.77\times10^{-7}$), its final-pass contrast ($p=7.63\times10^{-5}$), and the
protocol gap on 300 random items ($p$=3.1$\times$10$^{-6}$). \textbf{Not surviving it}: the
repair-loop pass gain ($p$=0.093 at $n$=100, $p$=0.0385 at $n$=300), the Llama protocol contrast
at $n$=300 ($p$=0.0352), the $\tau$-bench outcome difference ($p$=0.25), and the residual-gap
estimates ($p$=0.118, 0.125). We rest no claim on the latter group.

\textbf{The historical training verifier had two exploitable seams} (item~14): an unanchored
\texttt{passed} regex and a denominator omitting skipped/expected-failure tests. We audited the
historical logs; before the formal 7B control we also anchored the regex, counted all pytest
outcomes, and retained complete rollout text with hashes. Both P3 arms use that same repaired
verifier, so verifier changes cannot explain their contrast; the older learning curves remain
supported by an audit rather than a rerun.

Also constraining, and stated in full in the appendix: one task family and one tool (item~1); a
non-random 100-item set, whose 7B rung the random 300-item replication moves from 21 to 30
(item~2); a deliberately asymmetric taxonomy in which Mistral has no admissible FC comparison at
all (item~4); single-sample headline arms (item~5); scale stopping at 32B (item~7); sufficiency
tested by one repaired-FC seed and budget (item~12); a stratified validation sample with a 4000-character output
cap that creates false negatives while classifier imprecision creates false positives, leaving
the raw count's net bias unidentified (item~13); dependence on vLLM 0.27.1, verl 0.9.0
and publisher-mutable chat templates (items~15--16); and the instrumentation errata we disclose
rather than absorb --- seven arms recording \texttt{max\_tokens} they did not run at, Llama's
initiation rate corrected from 97 to 74, and one previously reported $p$=0.648 withdrawn
(item~17, Appendix~\ref{app:errata}).

\subsection*{What would change our conclusions}
\label{sec:future}

We name these so the claims are falsifiable rather than merely hedged.
\emph{Replications or a longer repaired-FC budget.} The completed single-variable control restores
16{,}844 tool interactions without moving held-out rescues. A consistent rescue gain over
additional seeds or a longer preregistered horizon would narrow the scope of this 150-step,
single-seed observation; continued nulls would strengthen it within, but not beyond, the tested
model/task setting.
\emph{A second family's scale ladder}, with a working test executor: if the undercount does not
grow with checkpoint size elsewhere, \S\ref{sec:scale}'s claim is a Qwen-plus-\texttt{hermes} fact
rather than a scaling phenomenon.
\emph{The same ladder on a non-code task.} The five Qwen2.5-Coder checkpoints run against
$\tau$-bench \texttt{retail} under the documented configuration --- models and interface fixed,
only the domain varied. If the emitted column still rises with size while the parsed column stays
at zero, the scale result is a property of the interface; if it does not, it is a property of
code-shaped payloads. This needs a rescue arm at the small sizes for the reason Granite did: a
checkpoint that simply never attempts a call produces an uninformative zero, not a censored one.
\emph{The same five-layer measurement on further benchmarks.} Done for BFCL~v4 and $\tau$-bench
\texttt{retail}, so the concern is not specific to our harness and now covers one interactive
suite; suites with much larger tool schemas, or whose scoring depends on argument-level matching,
remain untested.
